% Clase 6 --- Carga puntual y esfera centrada (§2.2).
% En toda la esfera E es radial (paralelo a n) y de módulo constante, así que
% Phi = E * 4 pi r^2 = Q/eps0, INDEPENDIENTE del radio: lo que se pierde en E
% (1/r^2) se gana en área (r^2).
\begin{center}
\begin{tikzpicture}[scale=1]

  % dos esferas concéntricas, para mostrar la independencia del radio
  \draw[figline] (0,0) circle (1.3);
  \draw[figaux] (0,0) circle (2.45);

  % carga
  \node[figpos, minimum size=8pt] at (0,0) {};
  \node[figlbl, below left=2pt] at (0,0) {$Q$};

  % campo radial: flechas largas en la esfera chica, cortas en la grande
  \foreach \a in {20,70,...,340} {
    \draw[figvec] (\a:1.3) -- (\a:2.05);
  }
  \foreach \a in {45,95,...,355} {
    \draw[figvec] (\a:2.45) -- (\a:2.85);
  }

  % radios, con los rótulos fuera de las flechas
  \draw[figvecaux] (0,0) -- (250:1.3);
  \node[figlblaux, anchor=north east] at (250:0.85) {$r$};
  \draw[figvecaux] (0,0) -- (300:2.45);
  \node[figlblaux, anchor=north west] at (300:1.55) {$r'$};

  % rótulos de lectura, fuera de todo
  \node[figlbl, figblue, anchor=west] at (3.15,1.4)
    {$\vec E \parallel \hat n$ \ en toda la esfera};
  \node[figlblaux, anchor=west] at (3.15,0.5) {$r$: \ $E$ mayor, menos área};
  \node[figlblaux, anchor=west] at (3.15,-0.15) {$r'$: \ $E$ menor, más área};

\end{tikzpicture}\\[0.5em]
{\small\itshape La esfera grande tiene un campo más débil pero más área, y en la
proporción exacta: $E$ cae como $1/r^{2}$ y el área crece como $r^{2}$. Por eso
el flujo \textbf{no depende del radio}. Todavía no es la ley de Gauss ---es sólo
una consecuencia de Coulomb---, pero ya tiene su forma.}
\end{center}
